\begin{center}
\refstepcounter{table}\label{tab:selection_function}
\begin{minipage}{\columnwidth}
\small\textbf{Table \thetable.} Gaia DR3 RVS selection-function
diagnostics for the catalogue.  The catalogue is selected from the
observed Gaia DR3 6D/RVS subset and is therefore not a volume-complete
Milky Way census.  $N_{\rm valid}$ is the number of sources for which
the GaiaUnlimited DR3-RVS implementation
\citep{CastroGinard2023} returned a finite selection-function weight at the source's $(G, BP{-}RP, l, b)$ cell;
the $N-N_{\rm valid}$ deficit is dominated by cells with zero parent
count.  $f_{n<10}$ is the fraction of $N_{\rm valid}$ sources whose
underlying GaiaUnlimited cell carries a parent count below 10 (the
regime in which the Beta$(k{+}1,\,n{-}k{+}1)$ posterior remains
competitive with its Beta$(1,1)$ prior: posterior standard error
$\gtrsim 0.16$ for $n<10$); $f_{\sum,n<10}$ is the share of
$\sum 1/S_{\rm RVS}$ contributed by those prior-dominated cells.  Tiers
are the adopted population-prior ones.  The
lower block restricts each tier to the inner-Galaxy quadrant
$l \in [330\degr,\,30\degr]$ where the catalogue concentrates
(Section~\ref{sec:sky}).  Inverse-selection weighted sums are reported
only as contextual observability diagnostics and are not used as
primary population counts.
\end{minipage}
\vspace{0.4ex}

\footnotesize
\setlength{\tabcolsep}{1pt}
\begin{tabular}{@{}lrrrrrr@{}}
\hline\hline
Sample & $N$ & $N_{\rm val}$ & $\Sigma 1/S$ & $\tilde S$ & $f_{<10}$ & $f_{\Sigma,<10}$ \\
\hline
\multicolumn{7}{l}{\textit{Full sample}} \\
Candidate pool          & 20{{,}}829 & 19{{,}}999 & 37{,}256.3 & 0.50 & --      & --      \\
Tier A                  & 173      & 172      & 273.5 & 0.67 & 84.3\%  & 89.6\%  \\
Tier B                  & 103      & 101      & 163.7 & 0.60 & 76.2\%  & 83.9\%  \\
Tier C                  & 345      & 343      & 584.7 & 0.50 & 84.5\%  & 90.0\%  \\
Tier A+B                & 276      & 273      & 437.2 & 0.67 & 81.3\%  & 87.5\%  \\
Tier A+B+C              & 621      & 616      & 1{,}021.9 & 0.67 & 83.1\%  & 88.9\%  \\
\hline
\multicolumn{7}{l}{\textit{Inner-Galaxy quadrant, $l\in[330\degr,\,30\degr]$}} \\
Tier A (inner)          & 17       & 17       & 22.1 & 0.90 & 52.9\% & 60.7\% \\
Tier A+B (inner)        & 32       & 31       & 43.5 & 0.83 & 54.8\% & 63.5\% \\
Tier A+B+C (inner)      & 88       & 87       & 138.6 & 0.67 & 67.8\% & 76.3\% \\
\hline
\end{tabular}
\end{center}
